\begin{abstract}
Automotive Ethernet (AE) is becoming the primary in-vehicle communication backbone for bandwidth-intensive applications such as camera streaming and ADAS, yet its security remains underexplored. Existing in-vehicle IDS predominantly rely on supervised learning or protocol-specific parsing, requiring labeled attack data and domain expertise difficult to obtain in real deployments.

We propose a protocol-agnostic, unsupervised IDS based on a dual-stream masked autoencoder trained exclusively on unlabeled normal traffic. The model processes 64-packet windows through a Transformer encoder over raw byte content and a residual 1D-CNN over inter-packet byte differences, fused via a delta-conditioned gated residual mechanism into a latent representation through attention pooling. Rather than the conventional sinusoidal scheme, the Transformer's positional encoding is conditioned on inter-packet capture timestamps, allowing it to jointly reason about content and arrival timing without protocol-specific parsing. Training uses packet-level masked reconstruction with a learnable mask token; at inference, anomalies are flagged by thresholding reconstruction error against the normal validation error distribution.

We evaluate on CarDS, a recently released real-vehicle AE dataset, using over 5.4 million normal packets for training and 27 attack files spanning 23 attack types, including network scanning, IP/MAC spoofing, and RTP stream injection. The proposed method achieves an F1-score of 0.9802, AUROC of 0.9976, and AUPRC of 0.9974, with near-perfect detection on scanning attacks (F1 up to 0.9993) and competitive performance on payload-substitution attacks. Ablation studies confirm that both the dual-stream design and time-aware encoding independently improve detection, with the latter substantially widening the threshold range over which performance remains stable.
\end{abstract}